\begin{proof}[Supplement to D.45: preserving multipliers and the original fibre]
The same obstruction also excludes a peripheral lift of the original
recorded \(\alpha_{33,R}\) to the full relative pro-\(2\) source
\(\mathcal P_2\) of D.34. Consequently the specific fibre
over \([\alpha_{33,R}]\) in (D.35.5a) is empty. This strengthens
the exclusion of the D.39--D.44 actions without excluding a
replacement action having their compact CM coordinate or
their prescribed CM pair.

Here are the full source maps and multiplier argument.
In the literal sphere graph of D.45, \(t^2\) is conjugate to
the free generator \(f_2\). Its eighth power has a nonzero
degree-eight mod-two Magnus term, while its sixteenth power
has no term below degree sixteen. Thus \(t\) has exact
order \(32\) in both \(B/D_9F\) and \(B/D_{10}F\).
The program \texttt{d\_cm\_rank2\_multiplier\_transport.py}
checks this directly by tuple-monomial Cauchy multiplication:
the first degrees of \(t^2,t^4,t^8,t^{16}\) are
\(1,2,4,8\), and \(t^{32}=1\) through degree nine.
Every \(\lambda_2\) preserving the recorded lower fixed
cusp therefore satisfies \(\lambda_2\equiv1\pmod{32}\)
and also fixes that cusp in the next quotient. There is
no additional degree-nine multiplier column.

The relative source has the full cusp exponent pair
\((\lambda_2,1)\) from D.34. Its order-three component
is one. On the retained \(Q_0\), \(s,r\) keep their
actual order-two and order-three conjugacy classes;
the power of \(r\) is one, and the power of \(s\)
is one because \(\lambda_2\) is odd. The image of
\(t\) in \(Q_0\) has order six, and
\(c_0=t^6\) lies in \(K^{(2)}\).
These are the exact conditions used by the
D.45 centralizer normalization.


We make explicit that every normalized representative in
\(\operatorname{Aut}_{Q_0,t}(\mathcal P_2)\) has the torsion
conjugators required here. Let \(H\) be the discrete preimage
of \(\langle s\rangle<Q_0\). Its nine-coset action has three
fixed \(s\)-cosets and no fixed \(r\)-cosets; the marked
Bass--Serre graph has nine edges, six order-two vertices and
three order-three vertices, with three nontrivial vertex
stabilizers, all of order two. Its graph has first Betti
number one. Thus \(H=C_2*C_2*C_2*\mathbb Z\), with each
finite factor a discrete \(B\)-conjugate of \(\langle s\rangle\).
The program \texttt{verify\_torsion\_bridge.py} records the actual
cosets, permutations and the three conjugating words.
Since \([H:K]=2\), taking cores in \(H\) of normal
subgroups of \(2\)-power index in \(K\) proves that the
induced pro-two topology on \(K\) is full. The resulting
map of completed extensions identifies the preimage of
\(\langle s\rangle\) in \(\mathcal P_2\) with
\(H^{(2)}=C_2\amalg_2C_2\amalg_2C_2\amalg_2\mathbb Z_2\),
with precisely the marking in (D.35.15).
Herfort--Ribes' torsion theorem
\cite[Theorem A$'$, pp.~157--159]{A_HerfortRibes1985}
therefore makes the image of \(s\) under any such
automorphism conjugate to a finite vertex generator,
hence to \(s\) in \(\mathcal P_2\).
The subgroups generated by \(r\) and its image are
complements to \(K_2\) over the same order-three
subgroup of \(Q_0\). Profinite coprime complement
conjugacy gives a \(K_2\)-conjugator; the fixed
\(Q_0\) image fixes the exponent as one.

The sphere projection is preserved as well. The preimage
\(\widetilde F\) of \(\ker(Q_0\to S_3)\) is invariant
under an automorphism inducing the identity on \(Q_0\).
Its maximal pro-two quotient is the marked \(F^{(2)}\):
every finite two-group quotient of \(F\) restricts to
one of \(K\), hence factors through \(\widetilde F\),
and the converse follows by restriction to the dense
image of \(F\). Its characteristic pro-two residual
is consequently preserved. The induced quotient of
\(\mathcal P_2\) is exactly \(\mathcal P_F\) in
(D.45.2), and the characteristic subgroups
\(D_9F^{(2)},D_{10}F^{(2)}\) are preserved there.
This proves the needed automorphism transfer to both
sphere quotients before applying their obstruction.

Now \(\operatorname{Int}(c_0^{258})\) is an actual
continuous inner automorphism of \(\mathcal P_2\),
with inverse \(\operatorname{Int}(c_0^{-258})\).
It induces the identity on \(Q_0\) and fixes \(t\)
as an element, not merely as a conjugacy class.
It preserves every peripheral class and commutes
with taking the indicated cusp powers. Its
reduction on \(R=T_3\) is the exact inner map
in (D.39.6). In that quotient \(c_0\) has order
sixteen, and \(258-146=112\) is divisible by
sixteen, proving the two displayed forms agree.
A relative lift of \(\alpha_{33,R}\), composed
on the left with this source inner map, would
therefore project to the excluded D.39 sphere
action. Its preserving multiplier is invisible
in the next sphere quotient by the order
calculation above. D.45's exhaustive conjugator
argument applies and gives a contradiction.

Finally consider the precise quotient in
(D.35.5a). By the exact restriction sequence
(D.35.5), an element of the stated outer fibre
has a normalized relative representative.
Its reduction differs from
\(\alpha_{33,R}\) by a power of
\(\operatorname{Int}(c_0)\).
The map
\(\langle c_0\rangle\simeq\mathbb Z_2
\to\langle c_0\bmod D_9K\rangle=C_{16}\)
is onto. Hence that finite discrepancy lifts
to an actual \(c_0^z\) in the source and can
be removed by its inverse inner automorphism.
This would give an exact relative lift of
\(\alpha_{33,R}\), already excluded.
It proves emptiness of that particular outer
fibre with both directions of the
normalization transport retained.
\end{proof}
